\begin{abstract}
\small
Quantum reservoir computing processes temporal data through fixed quantum
dynamics, but operation requires controlled forgetting. In QRC, engineered
dissipation is usually tuned through one rate; here we show that the
organization of environmental coupling is an architectural degree of freedom.
We define two coordinates: jump family, which specifies the coupled
system-operator combinations, and rate profile, which distributes assigned
coupling weight within a family. In paired driven-spin reservoirs, equal-phase
collective relaxation provides a short-term-memory advantage over
uniform local relaxation. The ordering recurs under separate activity-matching,
midpoint-gap-matching, operating-point-selection, long-washout, and
reset-encoding controls. Matched
rank-one interventions show that rotating the supported Kossakowski
direction changes accessible memory while Kossakowski rank, nonzero spectrum
and trace, sitewise diagonal weights, coefficient magnitudes, assigned operator
weight, and the processor remain fixed. Across the fixed-weight task map, no
continuously driven design leads every task. Demonstrating that coupling organization changes
readout-accessible memory beyond scalar-rate tuning is a significant
architecture-level advance: it establishes coupling organization as a
direct computational design variable.
\end{abstract}
